\documentclass[linenumbers,trackchanges,twocolumn]{aastex701}
\usepackage{lipsum}

%%%% TIKZ PREAMBLE 
% Network-diagram preamble: % Network-diagram preamble: % Network-diagram preamble: \input{architecture_preamble} before \begin{document}.
% Packages, colours, styles and the \mdfmap/\mdslabs/\mdimstack macros that
% architecture_figure.tex needs; thumbnails are read from images/thumbs.
%% --- network diagram (Universal 257) ---
\usepackage{tikz}
\usepackage{graphicx}
\providecommand{\mdthumbs}{images/thumbs}
\usetikzlibrary{arrows.meta,calc,positioning,backgrounds,fit}
\definecolor{mdLR}{HTML}{8A9BA8}     % feature maps, LR grid
\definecolor{mdHR}{HTML}{2E9E73}     % feature maps, SR grid
\definecolor{mdRes}{HTML}{E69F00}    % residual block
\definecolor{mdKnee}{HTML}{8E6BB8}   % learned output knee
\definecolor{mdLine}{HTML}{2B2B2B}
\tikzset{
  >={Stealth[length=4.6pt,width=3.8pt]},
  mdflow/.style={->,line width=0.75pt,draw=mdLine},
  mdop/.style={font=\fontsize{9}{10.5}\selectfont,align=center,inner sep=1.5pt},
  mdch/.style={font=\fontsize{8}{9.5}\selectfont,text=black!60,inner sep=1pt},
  mdsum/.style={circle,draw=mdLine,line width=0.75pt,fill=white,minimum size=12pt,inner sep=0pt},
  mdbox/.style={rounded corners=2pt,line width=0.7pt,align=center,font=\fontsize{9}{10.5}\selectfont,inner sep=4pt},
  mdpanel/.style={font=\fontsize{9}{10.5}\selectfont\bfseries,anchor=north west,inner sep=0pt},
}
% \mdfmap{name}{x}{y}{thickness}{height}{colour}: feature-map cuboid
\newcommand{\mdfmap}[6]{%
  \pgfmathsetmacro{\mdFx}{#2}\pgfmathsetmacro{\mdFy}{#3}%
  \pgfmathsetmacro{\mdFt}{#4}\pgfmathsetmacro{\mdFh}{#5}%
  \pgfmathsetmacro{\mdFdx}{0.30*\mdFh}\pgfmathsetmacro{\mdFdy}{0.22*\mdFh}%
  \begin{scope}[line width=0.5pt,draw=#6!50!black,line join=round]
    \filldraw[fill=#6!22] (\mdFx,\mdFy+0.5*\mdFh) -- ++(\mdFt,0) -- ++(\mdFdx,\mdFdy) -- ++(-\mdFt,0) -- cycle;
    \filldraw[fill=#6!62] (\mdFx+\mdFt,\mdFy-0.5*\mdFh) -- ++(0,\mdFh) -- ++(\mdFdx,\mdFdy) -- ++(0,-\mdFh) -- cycle;
    \filldraw[fill=#6!42] (\mdFx,\mdFy-0.5*\mdFh) rectangle ++(\mdFt,\mdFh);
  \end{scope}
  \coordinate (#1-w) at (\mdFx,\mdFy);
  \coordinate (#1-e) at (\mdFx+\mdFt+\mdFdx,\mdFy);
  \coordinate (#1-n) at (\mdFx+0.5*\mdFt+0.5*\mdFdx,\mdFy+0.5*\mdFh+\mdFdy);
  \coordinate (#1-s) at (\mdFx+0.5*\mdFt,\mdFy-0.5*\mdFh);
}
% \mdslabs{x}{y}{thickness}{height}{n}{colour}: n equal slabs (the 8 knees)
\newcommand{\mdslabs}[6]{%
  \pgfmathsetmacro{\mdSx}{#1}\pgfmathsetmacro{\mdSy}{#2}%
  \pgfmathsetmacro{\mdSt}{#3}\pgfmathsetmacro{\mdSh}{#4}%
  \pgfmathsetmacro{\mdSdx}{0.30*\mdSh}\pgfmathsetmacro{\mdSdy}{0.22*\mdSh}%
  \pgfmathtruncatemacro{\mdSn}{#5-1}%
  \foreach \i in {1,...,\mdSn}{%
    \pgfmathsetmacro{\mdXi}{\mdSx+\i*\mdSt/#5}%
    \draw[#6!50!black,line width=0.35pt]
      (\mdXi,\mdSy-0.5*\mdSh) -- (\mdXi,\mdSy+0.5*\mdSh) -- ++(\mdSdx,\mdSdy);
  }%
}
% \mdimstack{name}{x}{y}{size}{image}: colour image in front of three planes
\newcommand{\mdimstack}[5]{%
  \pgfmathsetmacro{\mdBx}{#2}\pgfmathsetmacro{\mdBy}{#3}\pgfmathsetmacro{\mdBs}{#4}%
  \foreach \k in {3,2,1}{%
    \filldraw[fill=white,draw=black!50,line width=0.4pt]
      ({\mdBx+0.07*\k},{\mdBy-0.5*\mdBs+0.07*\k}) rectangle ++(\mdBs,\mdBs);
  }%
  \node[inner sep=0pt,anchor=south west] at (\mdBx,\mdBy-0.5*\mdBs)
    {\includegraphics[width=\mdBs cm]{\mdthumbs/#5.png}};
  \draw[black!80,line width=0.4pt] (\mdBx,\mdBy-0.5*\mdBs) rectangle ++(\mdBs,\mdBs);
  \coordinate (#1-w) at (\mdBx,\mdBy);
  \coordinate (#1-e) at (\mdBx+\mdBs+0.21,\mdBy);
  \coordinate (#1-s) at (\mdBx+0.5*\mdBs,\mdBy-0.5*\mdBs);
}
\tikzset{
  mdline/.style={line width=0.75pt,draw=mdLine},
  mdfun/.style={rounded corners=6pt,draw=mdKnee!80!black,fill=mdKnee!16,line width=0.7pt,
    font=\fontsize{9}{10.5}\selectfont,align=center,inner xsep=4pt,inner ysep=3pt},
  mdleg/.style={font=\fontsize{8}{9.5}\selectfont,text=black!75,inner sep=1pt},
  mdnote/.style={font=\fontsize{8}{9.5}\selectfont\itshape,text=black!60,inner sep=1pt},
}
% Curve glyphs for the pixel-wise functions (small axes + the curve)
\newcommand{\mdgaxes}{\draw[black!35,line width=0.3pt] (-0.21,0) -- (0.21,0) (0,-0.19) -- (0,0.19);}
\newcommand{\mdgasinh}{\tikz[baseline=-0.6ex,sharp corners,line cap=round]{\mdgaxes
  \draw[mdKnee!70!black,line width=0.8pt,samples=41,domain=-0.2:0.2]
    plot (\x,{0.17*ln(30*\x+sqrt((30*\x)^2+1))/ln(6+sqrt(37))});}}
\newcommand{\mdgsinh}{\tikz[baseline=-0.6ex,sharp corners,line cap=round]{\mdgaxes
  \draw[mdKnee!70!black,line width=0.8pt,samples=41,domain=-0.2:0.2]
    plot (\x,{0.17*sinh(15*\x)/sinh(3)});}}
\newcommand{\mdgrelu}{\tikz[baseline=-0.6ex,sharp corners,line cap=round]{\mdgaxes
  \draw[mdKnee!70!black,line width=0.8pt,line join=round] (-0.2,0) -- (0,0) -- (0.17,0.17);}}
% \mdtwo{line 1}{line 2}: two centred lines, vertically centred on the math
% axis (AASTeX's tabular pads the box, so a plain \halign is used instead)
\newcommand{\mdtwo}[2]{$\vcenter{\halign{\hfil##\hfil\cr\strut#1\cr\strut#2\cr}}$}
 before \begin{document}.
% Packages, colours, styles and the \mdfmap/\mdslabs/\mdimstack macros that
% architecture_figure.tex needs; thumbnails are read from images/thumbs.
%% --- network diagram (Universal 257) ---
\usepackage{tikz}
\usepackage{graphicx}
\providecommand{\mdthumbs}{images/thumbs}
\usetikzlibrary{arrows.meta,calc,positioning,backgrounds,fit}
\definecolor{mdLR}{HTML}{8A9BA8}     % feature maps, LR grid
\definecolor{mdHR}{HTML}{2E9E73}     % feature maps, SR grid
\definecolor{mdRes}{HTML}{E69F00}    % residual block
\definecolor{mdKnee}{HTML}{8E6BB8}   % learned output knee
\definecolor{mdLine}{HTML}{2B2B2B}
\tikzset{
  >={Stealth[length=4.6pt,width=3.8pt]},
  mdflow/.style={->,line width=0.75pt,draw=mdLine},
  mdop/.style={font=\fontsize{9}{10.5}\selectfont,align=center,inner sep=1.5pt},
  mdch/.style={font=\fontsize{8}{9.5}\selectfont,text=black!60,inner sep=1pt},
  mdsum/.style={circle,draw=mdLine,line width=0.75pt,fill=white,minimum size=12pt,inner sep=0pt},
  mdbox/.style={rounded corners=2pt,line width=0.7pt,align=center,font=\fontsize{9}{10.5}\selectfont,inner sep=4pt},
  mdpanel/.style={font=\fontsize{9}{10.5}\selectfont\bfseries,anchor=north west,inner sep=0pt},
}
% \mdfmap{name}{x}{y}{thickness}{height}{colour}: feature-map cuboid
\newcommand{\mdfmap}[6]{%
  \pgfmathsetmacro{\mdFx}{#2}\pgfmathsetmacro{\mdFy}{#3}%
  \pgfmathsetmacro{\mdFt}{#4}\pgfmathsetmacro{\mdFh}{#5}%
  \pgfmathsetmacro{\mdFdx}{0.30*\mdFh}\pgfmathsetmacro{\mdFdy}{0.22*\mdFh}%
  \begin{scope}[line width=0.5pt,draw=#6!50!black,line join=round]
    \filldraw[fill=#6!22] (\mdFx,\mdFy+0.5*\mdFh) -- ++(\mdFt,0) -- ++(\mdFdx,\mdFdy) -- ++(-\mdFt,0) -- cycle;
    \filldraw[fill=#6!62] (\mdFx+\mdFt,\mdFy-0.5*\mdFh) -- ++(0,\mdFh) -- ++(\mdFdx,\mdFdy) -- ++(0,-\mdFh) -- cycle;
    \filldraw[fill=#6!42] (\mdFx,\mdFy-0.5*\mdFh) rectangle ++(\mdFt,\mdFh);
  \end{scope}
  \coordinate (#1-w) at (\mdFx,\mdFy);
  \coordinate (#1-e) at (\mdFx+\mdFt+\mdFdx,\mdFy);
  \coordinate (#1-n) at (\mdFx+0.5*\mdFt+0.5*\mdFdx,\mdFy+0.5*\mdFh+\mdFdy);
  \coordinate (#1-s) at (\mdFx+0.5*\mdFt,\mdFy-0.5*\mdFh);
}
% \mdslabs{x}{y}{thickness}{height}{n}{colour}: n equal slabs (the 8 knees)
\newcommand{\mdslabs}[6]{%
  \pgfmathsetmacro{\mdSx}{#1}\pgfmathsetmacro{\mdSy}{#2}%
  \pgfmathsetmacro{\mdSt}{#3}\pgfmathsetmacro{\mdSh}{#4}%
  \pgfmathsetmacro{\mdSdx}{0.30*\mdSh}\pgfmathsetmacro{\mdSdy}{0.22*\mdSh}%
  \pgfmathtruncatemacro{\mdSn}{#5-1}%
  \foreach \i in {1,...,\mdSn}{%
    \pgfmathsetmacro{\mdXi}{\mdSx+\i*\mdSt/#5}%
    \draw[#6!50!black,line width=0.35pt]
      (\mdXi,\mdSy-0.5*\mdSh) -- (\mdXi,\mdSy+0.5*\mdSh) -- ++(\mdSdx,\mdSdy);
  }%
}
% \mdimstack{name}{x}{y}{size}{image}: colour image in front of three planes
\newcommand{\mdimstack}[5]{%
  \pgfmathsetmacro{\mdBx}{#2}\pgfmathsetmacro{\mdBy}{#3}\pgfmathsetmacro{\mdBs}{#4}%
  \foreach \k in {3,2,1}{%
    \filldraw[fill=white,draw=black!50,line width=0.4pt]
      ({\mdBx+0.07*\k},{\mdBy-0.5*\mdBs+0.07*\k}) rectangle ++(\mdBs,\mdBs);
  }%
  \node[inner sep=0pt,anchor=south west] at (\mdBx,\mdBy-0.5*\mdBs)
    {\includegraphics[width=\mdBs cm]{\mdthumbs/#5.png}};
  \draw[black!80,line width=0.4pt] (\mdBx,\mdBy-0.5*\mdBs) rectangle ++(\mdBs,\mdBs);
  \coordinate (#1-w) at (\mdBx,\mdBy);
  \coordinate (#1-e) at (\mdBx+\mdBs+0.21,\mdBy);
  \coordinate (#1-s) at (\mdBx+0.5*\mdBs,\mdBy-0.5*\mdBs);
}
\tikzset{
  mdline/.style={line width=0.75pt,draw=mdLine},
  mdfun/.style={rounded corners=6pt,draw=mdKnee!80!black,fill=mdKnee!16,line width=0.7pt,
    font=\fontsize{9}{10.5}\selectfont,align=center,inner xsep=4pt,inner ysep=3pt},
  mdleg/.style={font=\fontsize{8}{9.5}\selectfont,text=black!75,inner sep=1pt},
  mdnote/.style={font=\fontsize{8}{9.5}\selectfont\itshape,text=black!60,inner sep=1pt},
}
% Curve glyphs for the pixel-wise functions (small axes + the curve)
\newcommand{\mdgaxes}{\draw[black!35,line width=0.3pt] (-0.21,0) -- (0.21,0) (0,-0.19) -- (0,0.19);}
\newcommand{\mdgasinh}{\tikz[baseline=-0.6ex,sharp corners,line cap=round]{\mdgaxes
  \draw[mdKnee!70!black,line width=0.8pt,samples=41,domain=-0.2:0.2]
    plot (\x,{0.17*ln(30*\x+sqrt((30*\x)^2+1))/ln(6+sqrt(37))});}}
\newcommand{\mdgsinh}{\tikz[baseline=-0.6ex,sharp corners,line cap=round]{\mdgaxes
  \draw[mdKnee!70!black,line width=0.8pt,samples=41,domain=-0.2:0.2]
    plot (\x,{0.17*sinh(15*\x)/sinh(3)});}}
\newcommand{\mdgrelu}{\tikz[baseline=-0.6ex,sharp corners,line cap=round]{\mdgaxes
  \draw[mdKnee!70!black,line width=0.8pt,line join=round] (-0.2,0) -- (0,0) -- (0.17,0.17);}}
% \mdtwo{line 1}{line 2}: two centred lines, vertically centred on the math
% axis (AASTeX's tabular pads the box, so a plain \halign is used instead)
\newcommand{\mdtwo}[2]{$\vcenter{\halign{\hfil##\hfil\cr\strut#1\cr\strut#2\cr}}$}
 before \begin{document}.
% Packages, colours, styles and the \mdfmap/\mdslabs/\mdimstack macros that
% architecture_figure.tex needs; thumbnails are read from images/thumbs.
%% --- network diagram (Universal 257) ---
\usepackage{tikz}
\usepackage{graphicx}
\providecommand{\mdthumbs}{images/thumbs}
\usetikzlibrary{arrows.meta,calc,positioning,backgrounds,fit}
\definecolor{mdLR}{HTML}{8A9BA8}     % feature maps, LR grid
\definecolor{mdHR}{HTML}{2E9E73}     % feature maps, SR grid
\definecolor{mdRes}{HTML}{E69F00}    % residual block
\definecolor{mdKnee}{HTML}{8E6BB8}   % learned output knee
\definecolor{mdLine}{HTML}{2B2B2B}
\tikzset{
  >={Stealth[length=4.6pt,width=3.8pt]},
  mdflow/.style={->,line width=0.75pt,draw=mdLine},
  mdop/.style={font=\fontsize{9}{10.5}\selectfont,align=center,inner sep=1.5pt},
  mdch/.style={font=\fontsize{8}{9.5}\selectfont,text=black!60,inner sep=1pt},
  mdsum/.style={circle,draw=mdLine,line width=0.75pt,fill=white,minimum size=12pt,inner sep=0pt},
  mdbox/.style={rounded corners=2pt,line width=0.7pt,align=center,font=\fontsize{9}{10.5}\selectfont,inner sep=4pt},
  mdpanel/.style={font=\fontsize{9}{10.5}\selectfont\bfseries,anchor=north west,inner sep=0pt},
}
% \mdfmap{name}{x}{y}{thickness}{height}{colour}: feature-map cuboid
\newcommand{\mdfmap}[6]{%
  \pgfmathsetmacro{\mdFx}{#2}\pgfmathsetmacro{\mdFy}{#3}%
  \pgfmathsetmacro{\mdFt}{#4}\pgfmathsetmacro{\mdFh}{#5}%
  \pgfmathsetmacro{\mdFdx}{0.30*\mdFh}\pgfmathsetmacro{\mdFdy}{0.22*\mdFh}%
  \begin{scope}[line width=0.5pt,draw=#6!50!black,line join=round]
    \filldraw[fill=#6!22] (\mdFx,\mdFy+0.5*\mdFh) -- ++(\mdFt,0) -- ++(\mdFdx,\mdFdy) -- ++(-\mdFt,0) -- cycle;
    \filldraw[fill=#6!62] (\mdFx+\mdFt,\mdFy-0.5*\mdFh) -- ++(0,\mdFh) -- ++(\mdFdx,\mdFdy) -- ++(0,-\mdFh) -- cycle;
    \filldraw[fill=#6!42] (\mdFx,\mdFy-0.5*\mdFh) rectangle ++(\mdFt,\mdFh);
  \end{scope}
  \coordinate (#1-w) at (\mdFx,\mdFy);
  \coordinate (#1-e) at (\mdFx+\mdFt+\mdFdx,\mdFy);
  \coordinate (#1-n) at (\mdFx+0.5*\mdFt+0.5*\mdFdx,\mdFy+0.5*\mdFh+\mdFdy);
  \coordinate (#1-s) at (\mdFx+0.5*\mdFt,\mdFy-0.5*\mdFh);
}
% \mdslabs{x}{y}{thickness}{height}{n}{colour}: n equal slabs (the 8 knees)
\newcommand{\mdslabs}[6]{%
  \pgfmathsetmacro{\mdSx}{#1}\pgfmathsetmacro{\mdSy}{#2}%
  \pgfmathsetmacro{\mdSt}{#3}\pgfmathsetmacro{\mdSh}{#4}%
  \pgfmathsetmacro{\mdSdx}{0.30*\mdSh}\pgfmathsetmacro{\mdSdy}{0.22*\mdSh}%
  \pgfmathtruncatemacro{\mdSn}{#5-1}%
  \foreach \i in {1,...,\mdSn}{%
    \pgfmathsetmacro{\mdXi}{\mdSx+\i*\mdSt/#5}%
    \draw[#6!50!black,line width=0.35pt]
      (\mdXi,\mdSy-0.5*\mdSh) -- (\mdXi,\mdSy+0.5*\mdSh) -- ++(\mdSdx,\mdSdy);
  }%
}
% \mdimstack{name}{x}{y}{size}{image}: colour image in front of three planes
\newcommand{\mdimstack}[5]{%
  \pgfmathsetmacro{\mdBx}{#2}\pgfmathsetmacro{\mdBy}{#3}\pgfmathsetmacro{\mdBs}{#4}%
  \foreach \k in {3,2,1}{%
    \filldraw[fill=white,draw=black!50,line width=0.4pt]
      ({\mdBx+0.07*\k},{\mdBy-0.5*\mdBs+0.07*\k}) rectangle ++(\mdBs,\mdBs);
  }%
  \node[inner sep=0pt,anchor=south west] at (\mdBx,\mdBy-0.5*\mdBs)
    {\includegraphics[width=\mdBs cm]{\mdthumbs/#5.png}};
  \draw[black!80,line width=0.4pt] (\mdBx,\mdBy-0.5*\mdBs) rectangle ++(\mdBs,\mdBs);
  \coordinate (#1-w) at (\mdBx,\mdBy);
  \coordinate (#1-e) at (\mdBx+\mdBs+0.21,\mdBy);
  \coordinate (#1-s) at (\mdBx+0.5*\mdBs,\mdBy-0.5*\mdBs);
}
\tikzset{
  mdline/.style={line width=0.75pt,draw=mdLine},
  mdfun/.style={rounded corners=6pt,draw=mdKnee!80!black,fill=mdKnee!16,line width=0.7pt,
    font=\fontsize{9}{10.5}\selectfont,align=center,inner xsep=4pt,inner ysep=3pt},
  mdleg/.style={font=\fontsize{8}{9.5}\selectfont,text=black!75,inner sep=1pt},
  mdnote/.style={font=\fontsize{8}{9.5}\selectfont\itshape,text=black!60,inner sep=1pt},
}
% Curve glyphs for the pixel-wise functions (small axes + the curve)
\newcommand{\mdgaxes}{\draw[black!35,line width=0.3pt] (-0.21,0) -- (0.21,0) (0,-0.19) -- (0,0.19);}
\newcommand{\mdgasinh}{\tikz[baseline=-0.6ex,sharp corners,line cap=round]{\mdgaxes
  \draw[mdKnee!70!black,line width=0.8pt,samples=41,domain=-0.2:0.2]
    plot (\x,{0.17*ln(30*\x+sqrt((30*\x)^2+1))/ln(6+sqrt(37))});}}
\newcommand{\mdgsinh}{\tikz[baseline=-0.6ex,sharp corners,line cap=round]{\mdgaxes
  \draw[mdKnee!70!black,line width=0.8pt,samples=41,domain=-0.2:0.2]
    plot (\x,{0.17*sinh(15*\x)/sinh(3)});}}
\newcommand{\mdgrelu}{\tikz[baseline=-0.6ex,sharp corners,line cap=round]{\mdgaxes
  \draw[mdKnee!70!black,line width=0.8pt,line join=round] (-0.2,0) -- (0,0) -- (0.17,0.17);}}
% \mdtwo{line 1}{line 2}: two centred lines, vertically centred on the math
% axis (AASTeX's tabular pads the box, so a plain \halign is used instead)
\newcommand{\mdtwo}[2]{$\vcenter{\halign{\hfil##\hfil\cr\strut#1\cr\strut#2\cr}}$}


%%%% REVIEW COMMENTS (todonotes)
%% \commentstrue  -> notes show as coloured boxes + a list at the top
%% \commentsfalse -> every note disappears; the PDF is the clean paper
\newif\ifcomments \commentstrue
\ifcomments
  \usepackage[colorinlistoftodos,textsize=footnotesize]{todonotes}
\else
  \usepackage[disable]{todonotes}
\fi
\setuptodonotes{inline,bordercolor=black!40}
%% \ab{date}{text}: one comment from Aleksandr (as on Overleaf)
\newcommand{\ab}[2]{\todo[color=orange!25,author={Aleksandr, #1}]{#2}}
%% \board{text}: transcribed whiteboard sketch
\newcommand{\board}[1]{\todo[color=violet!15,author={Whiteboard},caption={Whiteboard: intro outline}]{#1}}
%% \task{text}: open task from the to-do list
\newcommand{\task}[1]{\todo[color=cyan!20,author={Task}]{#1}}
%% \opt{paragraph}{n}: label one of the alternative versions of a paragraph (hidden with \commentsfalse)
\newcommand{\opt}[2]{\par\medskip\noindent\ifcomments{\color{teal!70!black}\sffamily\bfseries\footnotesize #1 --- option #2}\par\noindent\fi}


\begin{document}

\title{\lipsum[1]}

\author[orcid=0000-0000-0000-0001,sname='North America']{Tundra North America}
\affiliation{Harvard COllege}
\email[show]{abelotserkovtsev@college.harvard.edu}  


%% Keywords should appear after the \end{abstract} command. 
%% The AAS Journals now uses Unified Astronomy Thesaurus (UAT) concepts:
%% https://astrothesaurus.org
%% You will be asked to selected these concepts during the submission process
%% but this old "keyword" functionality is maintained in case authors want
%% to include these concepts in their preprints.
%%
%% You can use the \uat command to link your UAT concepts back its source.
\keywords{\uat{Galaxies}{573} --- \uat{Cosmology}{343} --- \uat{High Energy astrophysics}{739} --- \uat{Interstellar medium}{847} --- \uat{Stellar astronomy}{1583} --- \uat{Solar physics}{1476}}

\ifcomments \listoftodos[Review comments] \fi

\task{Intro + presentation, alternating: presentation $\to$ introduction $\to$ presentation $\to$ introduction. A lot more focus on the introduction than in the last presentation; draft the slides.}

\section{Introduction}

\board{Intro outline (whiteboard):
\begin{itemize}
  \item \textbf{Wide field} optical imaging surveys: \textbf{always} from the ground. Can't discover cool things without looking at large areas!
  \item vs.\ space-based imaging with a tiny field of view (e.g.\ Hubble, JWST): expensive!
  \item $\to$ introduce SR studies on ground imaging (SDSS, HSC, DECaLS, CFHT, SkyMapper\ldots). SDSS and DECaLS: starting to find gravitational lenses (that Tania likes!)
  \item $\to$ introduce Liam's work: from the ground, but RADIO, which also has a stable PSF.
  \item $\to$ EUCLID! \textbf{First time}: large-area optical imaging \emph{and} a stable PSF.
  \item $\to$ also, Roman is upcoming.
\end{itemize}}

\ab{11 Nov, 12:00 pm}{Paragraph 1: what field --- optical imaging surveys.}
\ab{11 Nov, 12:00 pm}{History lesson.}
\ab{11 Nov, 12:00 pm}{Set up the fact that it's been for a while. BUT ALWAYS from the ground.}
\ab{11 Nov, 12:01 pm}{vs.\ space-based imaging has always had a very, very tiny FOV.}
\ab{11 Nov, 12:02 pm}{Introduce SR studies on ground imaging.}
\ab{11 Nov, 12:03 pm}{Introduce Liam's work.}

\ab{11 Nov, 12:05 pm}{RADIO has a stable PSF. We don't!}
\ab{11 Nov, 12:06 pm}{First time we have large-area optical imaging + a stable PSF.}
\ab{11 Nov, 12:06 pm}{Euclid is here and Roman is upcoming!!! Should be doing this work now!}
\ab{11 Nov, 12:06 pm}{Tomorrow is Euclid DR1!!!}

\ab{11 Nov, 11:56 am}{We should add more pretext to this.}
\ab{11 Nov, 11:56 am}{Biggest selling point: the first time for optical astronomy that we have wide-field space-based imaging --- the key selling point.}
\ab{11 Nov, 11:56 am}{The first time we could do this.}
\ab{11 Nov, 11:57 am}{No one before could have done this.}

\task{17 Nov draft: seven paragraphs in the whiteboard order (ground surveys $\to$ space is small $\to$ super-resolution from the ground $\to$ radio/POLISH $\to$ Euclid is the first $\to$ lenses $\to$ this work), three options each. Option 1 stays closest to the 11 Nov text; option 2 tells the history; option 3 is the shortest. Pick one per paragraph (or mix sentences) and delete the rest. New references are marked in references.bib; check them against ADS before submission.}

%% ---------------------------------------------------------------- P1
\opt{P1: wide-field surveys, always from the ground}{1}
Wide-field imaging surveys are how astronomy finds what is rare and measures what is common.
Rare objects, such as strong gravitational lenses, and the statistics of galaxy populations only emerge when large areas of the sky are observed.
For decades these surveys have been carried out from the ground, from the Sloan Digital Sky Survey \citep{York2000} to the Dark Energy Survey \citep{DES2016}, the DESI Legacy Imaging Surveys \citep{Dey2019} and the Hyper Suprime-Cam survey \citep{Aihara2018}, with the Rubin Observatory's LSST now starting \citep{Ivezic2019}.
Each generation went wider and deeper, but all of them observe through the atmosphere, which blurs every image to about $0.6''$ at the best sites and to $1''$ or more at most others.
Wide-field optical imaging has therefore always been limited in resolution by the atmosphere rather than by the telescope.

\opt{P1: wide-field surveys, always from the ground}{2}
Large optical surveys have a long history, and all of it has been written from the ground.
Since the photographic sky surveys of the mid-twentieth century, mapping large areas has been the way to find rare objects and to measure galaxy populations with enough statistics.
Digital surveys turned this into a precision science: SDSS set out to image a quarter of the sky \citep{York2000}, and the Dark Energy Survey \citep{DES2016}, the DESI Legacy Imaging Surveys \citep{Dey2019} and Hyper Suprime-Cam \citep{Aihara2018} went deeper, with LSST now beginning \citep{Ivezic2019}.
Yet every one of these surveys sees the sky through turbulent air, which spreads the light of each point source over about $0.6''$ at the best sites and over $1''$ or more at typical ones.
For seventy years, wide meant ground-based, and ground-based meant blurred by the atmosphere.

\opt{P1: wide-field surveys, always from the ground}{3}
Wide-field optical surveys find rare objects and measure galaxy populations, and until recently all of them were made from the ground \citep{York2000,DES2016,Dey2019,Aihara2018,Ivezic2019}.
Their resolution is set by the atmosphere, which blurs images to about $0.6''$ at the best sites and $1''$ or more elsewhere, not by the telescopes.

%% ---------------------------------------------------------------- P2
\opt{P2: space imaging is sharp but small}{1}
Space telescopes escape the atmosphere, but so far only over small fields.
This matters because the sharpest images of galaxies have come only from space, and only for small patches of the sky.
The \emph{Hubble Space Telescope} resolves structure on scales of $0.1''$, but its cameras cover about ten square arcminutes at a time, and its largest contiguous survey, COSMOS, covers about $2\,\mathrm{deg}^2$ \citep{Scoville2007COSMOS}.
JWST is sharper still, with a field of view of similar size.
Observing time on both telescopes is scarce, so they mostly follow up targets that wide ground-based surveys have already found.
Until now, astronomers have had to choose between images that are wide and images that are sharp.

\opt{P2: space imaging is sharp but small}{2}
The alternative to the ground has been space, at the price of area.
Hubble and JWST deliver images several times sharper than any ground-based survey, which is why they are used to study the structure of galaxies and to confirm lens candidates.
Their fields of view, however, are about ten square arcminutes, so even the largest contiguous Hubble survey, COSMOS, covers only about $2\,\mathrm{deg}^2$ \citep{Scoville2007COSMOS}, less than a thousandth of the area of a wide survey.
Time on these telescopes is oversubscribed, and covering thousands of square degrees with them is out of reach.
Space imaging has thus been a tool for following up a few objects, not for finding them.

\opt{P2: space imaging is sharp but small}{3}
Space telescopes avoid the atmosphere, but Hubble and JWST image only about ten square arcminutes at a time; the largest contiguous Hubble survey, COSMOS, covers about $2\,\mathrm{deg}^2$ \citep{Scoville2007COSMOS}.
Sharp images have been available for small fields and wide images only at ground-based resolution.

%% ---------------------------------------------------------------- P3
\opt{P3: super-resolution of ground-based images}{1}
Several groups have tried to close this gap by sharpening ground-based images with neural networks.
If such a network were reliable, every existing survey would gain resolution without new observations.
Recent work maps SDSS and DECaLS images to the quality of the deeper and sharper Hyper Suprime-Cam data \citep{Luo2025}, builds paired data sets of galaxy images for astronomical super-resolution \citep{Miao2024}, and deconvolves ground-based images towards space-like resolution \citep{Akhaury2024}.
All of these face the same difficulty: the atmospheric point-spread function (PSF) changes from exposure to exposure and across the field, so the blur that the network must undo is never known exactly.
A network that has learned one set of conditions can invent structure or bias measurements under another \citep{Akhaury2024}, and there is no wide, sharp reference to check it against.
Super-resolution of ground-based images is therefore possible, but hard to trust over a whole survey.

\opt{P3: super-resolution of ground-based images}{2}
Machine learning has made super-resolution of ground-based images an active field.
Networks now translate SDSS and DECaLS images towards Hyper Suprime-Cam quality \citep{Luo2025}, are trained and compared on shared galaxy data sets \citep{Miao2024}, and deconvolve ground-based observations towards a resolution of $0.1''$ \citep{Akhaury2024}.
These results are encouraging, but they rest on an unstable foundation.
The atmospheric point-spread function (PSF) varies with time, position and wavelength, so the training pairs can only approximate the blur in any given image, and the network's output can carry biases that are hard to detect \citep{Akhaury2024}.
From the ground, the main limit on super-resolution is not the network but the unknown PSF.

\opt{P3: super-resolution of ground-based images}{3}
Neural networks have been trained to sharpen ground-based images \citep{Luo2025,Miao2024,Akhaury2024}, but the atmospheric point-spread function (PSF) changes from exposure to exposure, so the blur they must undo is never known exactly and their outputs can be biased \citep{Akhaury2024}.

%% ---------------------------------------------------------------- P4
\opt{P4: radio interferometry and POLISH}{1}
Radio astronomy shows what becomes possible when the PSF is known.
An interferometer's PSF is set by the positions of its antennas, so it is known precisely for every observation, even from the ground.
\citet{Connor2022} used this to build POLISH, a feed-forward residual convolutional network that is trained on simulated skies and deconvolves and super-resolves radio images in a single pass, fast enough for the data rate of the DSA-2000.
On simulated observations POLISH outperformed the standard CLEAN algorithm and recovered sources smaller than the beam.
\citet{Wu2026} extended it to POLISH++, which processes wide fields in patches and compresses their dynamic range with an $\mathrm{arcsinh}$ stretch; combined with a lens classifier, it recovered strong-lens arcs at separations close to the size of the PSF.
With a stable, known PSF, a network trained only on simulations can reconstruct real survey images at scale.

\opt{P4: radio interferometry and POLISH}{2}
Radio interferometry faced the same need for fast, survey-scale reconstruction, but from a better starting point.
Next-generation arrays such as the DSA-2000 will produce data faster than iterative algorithms such as CLEAN can deconvolve it \citep{Connor2022}.
Their PSF, however, follows directly from the array geometry and is known for every observation.
\citet{Connor2022} introduced POLISH, a residual convolutional network trained on simulated sky--image pairs that deconvolves and super-resolves in one pass and beat CLEAN on simulated data.
POLISH++ \citep{Wu2026} added patch-wise processing for wide fields and an $\mathrm{arcsinh}$ stretch for high dynamic range, and recovered strong-lens morphology near the PSF scale.
These works established that simulation-trained reconstruction is accurate and fast when the PSF is stable.

\opt{P4: radio interferometry and POLISH}{3}
In radio astronomy the PSF follows from the array geometry and is known for every observation, and simulation-trained networks already reconstruct survey images: POLISH deconvolves and super-resolves in a single pass \citep{Connor2022}, and POLISH++ extends it to wide fields with high dynamic range and recovers strong-lens arcs near the PSF scale \citep{Wu2026}.

%% ---------------------------------------------------------------- P5
\opt{P5: Euclid is the first wide survey with a stable PSF}{1}
Euclid is the first optical and near-infrared survey to combine both conditions: a wide field and a stable, known PSF.
This is the combination that ground-based surveys and space telescopes have each had only half of.
Euclid observes from space in one visible and three near-infrared bands, and its PSF is stable and accurately modelled because its weak-lensing measurements require it \citep{Mellier2025}.
At the same time, its Wide Survey will cover about $14{,}000\,\mathrm{deg}^2$ \citep{Mellier2025}, several thousand times the area of the largest Hubble survey.
Its first large release, DR1-Foundation, came out in November 2026 with about $1900\,\mathrm{deg}^2$\footnote{\url{https://www.cosmos.esa.int/web/euclid/euclid-dr1}}, and the \emph{Nancy Grace Roman Space Telescope} \citep{Spergel2015}, launched in August 2026, will add more than $5000\,\mathrm{deg}^2$ of space-based near-infrared imaging\footnote{\url{https://www.nasa.gov/missions/roman-space-telescope/core-survey-by-nasas-roman-mission-will-unveil-universes-dark-side/}}.
For the first time, the approach that works for radio surveys can be applied to wide optical imaging, and the data to do it are arriving now.

\opt{P5: Euclid is the first wide survey with a stable PSF}{2}
This changed with \emph{Euclid}.
For the first time, a telescope images a large fraction of the sky from space, with the stable PSF that super-resolution needs.
Euclid's weak-lensing goals demand a sharp PSF that is stable and accurately modelled across the focal plane \citep{Mellier2025}, and its Wide Survey will cover about $14{,}000\,\mathrm{deg}^2$ in one visible ($I_{\rm E}$) and three near-infrared ($Y_{\rm E}$, $J_{\rm E}$, $H_{\rm E}$) bands \citep{Mellier2025}.
The first $1900\,\mathrm{deg}^2$ were released as DR1-Foundation in November 2026\footnote{\url{https://www.cosmos.esa.int/web/euclid/euclid-dr1}}, and the Roman Space Telescope \citep{Spergel2015}, launched in August 2026, will survey more than $5000\,\mathrm{deg}^2$ in the near-infrared\footnote{\url{https://www.nasa.gov/missions/roman-space-telescope/core-survey-by-nasas-roman-mission-will-unveil-universes-dark-side/}}.
No earlier data set could have supported survey-scale super-resolution of optical images; Euclid can, and now is the time to build it.

\opt{P5: Euclid is the first wide survey with a stable PSF}{3}
\emph{Euclid} is the first wide optical and near-infrared survey made from space: its PSF is stable and accurately modelled \citep{Mellier2025}, and its Wide Survey will cover about $14{,}000\,\mathrm{deg}^2$, of which $1900\,\mathrm{deg}^2$ were released in November 2026\footnote{\url{https://www.cosmos.esa.int/web/euclid/euclid-dr1}}.
With Roman \citep{Spergel2015}, launched in August 2026, adding more than $5000\,\mathrm{deg}^2$\footnote{\url{https://www.nasa.gov/missions/roman-space-telescope/core-survey-by-nasas-roman-mission-will-unveil-universes-dark-side/}}, survey-scale super-resolution of optical imaging is possible for the first time.

%% ---------------------------------------------------------------- P6
\opt{P6: strong lenses need the extra resolution}{1}
Strong gravitational lenses are one of the science cases that gain most from sharper wide-field images.
A strong lens forms when a massive foreground galaxy or cluster bends the light of a more distant source into multiple images or arcs \citep{Treu2010}.
Lenses magnify galaxies that would otherwise be too faint to study and measure the mass, including the dark matter, of the deflector and the geometry of the Universe \citep{Treu2010}.
They are rare, and ground-based surveys have found thousands of candidates \citep{Jacobs2019,Huang2020}, mostly with large Einstein radii, because a small arc blends into the light of its deflector.
Systems with small Einstein radii, often with lower-mass deflectors, are therefore underrepresented \citep{AcevedoBarroso2025}.
Euclid is expected to find more than $10^5$ lenses \citep{Collett2015_ApJ81120}, and its first $63\,\mathrm{deg}^2$ already yielded about 500 candidates \citep{Walmsley2025}; sharper images would extend these searches to the smallest and most blended systems.

\opt{P6: strong lenses need the extra resolution}{2}
The clearest case for higher resolution over wide areas is the search for strong gravitational lenses.
These rare systems, in which a foreground galaxy splits the light of a background source into arcs or multiple images, measure dark matter in the deflector and magnify sources too faint to see otherwise \citep{Treu2010}.
Finding them requires separating the arc from the deflector, which becomes impossible when the Einstein radius approaches the PSF size, so compact systems are missed and the known population is biased towards massive deflectors \citep{AcevedoBarroso2025}.
Euclid's first $63\,\mathrm{deg}^2$ already doubled the number of lens candidates found in space-based imaging \citep{Walmsley2025}, and the full survey is expected to contain more than $10^5$ lenses \citep{Collett2015_ApJ81120}.
Its four bands help further, because the source and the deflector usually differ in colour.
Resolution gained beyond Euclid's native sampling translates directly into lenses found at smaller Einstein radii.

\opt{P6: strong lenses need the extra resolution}{3}
Strong gravitational lenses are a direct beneficiary: they are rare, they measure dark matter and magnify faint sources \citep{Treu2010}, and those with small Einstein radii are missed when the arc blends with the deflector \citep{AcevedoBarroso2025}.
Euclid is expected to contain more than $10^5$ of them \citep{Collett2015_ApJ81120,Walmsley2025}.

%% ---------------------------------------------------------------- P7
\opt{P7: this work}{1}
In this work we adapt the POLISH framework to Euclid's visible and near-infrared imaging.
The model takes the four Euclid bands and reconstructs each field at twice the sampling, from $0.1''$ to $0.05''$ per pixel, while removing noise and the PSF.
Real Euclid images have no high-resolution counterparts, so we train the network on synthetic skies with a known truth.
The synthetic skies follow source populations measured in Euclid Q1 and take galaxy morphologies and colours from the TNG50-SKIRT Atlas \citep{Baes2024TNG50SKIRT}; each is turned into a Euclid-like observation with measured PSFs and models of the noise and detector artifacts.
We validate the model on synthetic images with known truth and then apply it to real Euclid images, comparing the reconstructions with JWST observations of the same galaxies.

\opt{P7: this work}{2}
Here we build the first step of such a pipeline for Euclid.
We adapt POLISH to the four Euclid bands, reconstructing each field on a $0.05''$ grid, twice as fine as the $0.1''$ input, while removing noise and the PSF in a single pass that is fast enough for the full survey.
Because no sharper images of the whole Euclid sky exist, the network learns from synthetic skies with known truth: galaxies from the TNG50-SKIRT Atlas \citep{Baes2024TNG50SKIRT} and stars, placed according to the populations measured in Euclid Q1, observed through Euclid's measured PSFs, noise and artifacts.
We test the model on synthetic images, where the truth is known, and on real Euclid Q1 images, where JWST observations of the same field provide a sharper view.
Section~\ref{sec:data} describes the data, the Methods describe the training set and the model, and the Results follow.

\opt{P7: this work}{3}
We adapt POLISH to Euclid: a network trained on synthetic skies built from Euclid Q1 populations and TNG50-SKIRT galaxies \citep{Baes2024TNG50SKIRT} reconstructs the four Euclid bands at $0.05''$ per pixel, and we test it on synthetic data and against JWST imaging of the same galaxies.

\section{Data}\label{sec:data}
We use three data sets:\begin{itemize}
    \item Euclid Q1 provides the images that we reconstruct and the measurements that calibrate our synthetic training data.
    \item The TNG50-SKIRT Atlas provides high-resolution images of synthetic realistic galaxy morphologies in Euclid bands, for the use in training data.
    \item NEXUS JWST imaging provides an independent, higher-resolution view of part of a Q1 field, which we use to assess the quality of reconstructions of real data.
\end{itemize}
Table~\ref{tab:imaging} summarizes the imaging characteristics of the Euclid and NEXUS data.

\begin{deluxetable*}{lcccccc}
\tablecaption{Imaging data used in this work.\label{tab:imaging}}
\tablehead{
\colhead{Band} & \colhead{Wavelength} & \colhead{Native pixel} & \colhead{Delivered pixel} & \colhead{Native PSF FWHM} & \colhead{Measured PSF FWHM} & \colhead{$5\sigma$ depth} \\
\colhead{} & \colhead{(nm)} & \colhead{(\arcsec)} & \colhead{(\arcsec)} & \colhead{(\arcsec)} & \colhead{(\arcsec)} & \colhead{(AB mag)}
}
\startdata
Euclid $I_{\rm E}$ & 530--920       & 0.10 & 0.10 & 0.13 & 0.147 & 26.7 \\
Euclid $Y_{\rm E}$ & 949.6--1212.3  & 0.30 & 0.10 & 0.33 & 0.363 & 24.4 \\
Euclid $J_{\rm E}$ & 1167.6--1567.0 & 0.30 & 0.10 & 0.35 & 0.379 & 24.4 \\
Euclid $H_{\rm E}$ & 1521.5--2021.4 & 0.30 & 0.10 & 0.35 & 0.395 & 24.4 \\
NEXUS F200W        & 1755--2227     & 0.031& 0.03 & 0.066& ---   & 27.7 \\
  \enddata
\tablecomments{Euclid bandpasses and wide-survey point-source depths from \citet{Mellier2025}. Euclid PSF FWHMs are medians measured from stars on the Q1 MER mosaics using Photutils. The NEXUS depth is single-epoch Deep-tier $5\sigma$ point-source depth \citep{Shen2024NEXUS}.}
\end{deluxetable*}

\subsection{Euclid Q1}\label{sec:data_euclid}
The Euclid data we use comes from the Euclid Quick Data Release Q1 \citep{aussel2025}, which covers $63.1\,\mathrm{deg}^2$ of the three Euclid Deep Fields (North, South, and Fornax) at the depth of a single Euclid Wide Survey visit.
Q1 is processed with the same pipeline that will be used for the Wide Survey, so training-inference pipeline on it transfers immediately to later releases.
VIS instrument provides imaging in the broad $I_{\rm E}$ band on $0.1''$ pixel grid and NISP images in $Y_{\rm E}$, $J_{\rm E}$, and $H_{\rm E}$ on $0.3''$ pixel grid \citep{Mellier2025}.

We use products of the MER processing function \citep{romelli2025} that subtracts background, removes part of the noise and artifacts, and resamples all four bands onto a common $0.1''\,\mathrm{pixel}^{-1}$ grid.
The bands nevertheless keep different resolutions: the median PSF FWHM in the mosaics is $0.15''$ in $I_{\rm E}$ and $0.36''$--$0.40''$ in the NISP bands (Table~\ref{tab:imaging}).
These values are $10$--$15\%$ larger than the PSFs of the individual calibrated frames because stacking and resampling broaden the PSF \citep{romelli2025}.
Because we are working only with the delivered mosaics, every PSF from now on will refer to the measured MER-level PSFs.

We also use the MER RMS maps that accompany each mosaic to model noise levels for synthetic training data; and the MER multi-band catalogue with the Q1 photometric-redshift products to build empirical distributions of galaxies' parameters for our synthetic training data.
Bright-star photometry is supplemented with \textit{Gaia} DR3 \citep{GaiaDR3}, as Euclid filters quickly become oversaturated by stars of  VIS AB magnitude below $16.5$.
All Euclid products were retrieved from the ESA Euclid Science Archive and IRSA.

\subsection{TNG50-SKIRT Atlas}\label{sec:data_tng}
Euclid does not resolve galaxy morphology at the scales the network must reconstruct, so we take morphologies from simulation.
TNG50-1 \citep{Pillepich2019,Nelson2019} is the highest-resolution run of the IllustrisTNG magnetohydrodynamical cosmological simulations.
The TNG50-SKIRT Atlas \citep{Baes2024TNG50SKIRT}, exteneded by \citet{Kovacic2025} to the four Euclid bands, post-processes every TNG50 galaxy at $z=0$ with stellar mass between $10^{9.8}$ and $10^{12}\,M_\odot$ ($1154$ galaxies, $869$ star-forming and $285$ quiescent) using the SKIRT radiative-transfer code.
SKIRT accounts for multiple stellar populations, absorption and scattering by interstellar dust, and attenuation in star-forming regions.

Each galaxy is imaged from five observer positions, and each image has $100\,\mathrm{pc}$ pixels over a $160\,\mathrm{kpc}$ field ($1600\times1600$ pixels).
The images contain no instrumental PSF and no noise, which is ideal for our use case, that requires high-resolution galaxy morphologies \citep{Kovacic2025}.
We use the resultant $I_{\rm E}Y_{\rm E}J_{\rm E}H_{\rm E}$ images in all five projections for $1140$ of the $1154$ galaxies.

These TNG50 galaxies are used to generate training data with high-resolution galaxies in them: we scale, rotate and change the brightnesses of these galaxies while preserving morphology (matched by star-formation rate) to fit the distribution of Euclid-measured galaxies.

\subsection{JWST NEXUS}\label{sec:data_nexus}
As a high-resolution reference we use JWST/NIRCam F200W imaging from NEXUS
\citep{Shen2024NEXUS}, whose field lies entirely within the Euclid Deep Field North and is therefore covered by Q1. 
We use the Deep-tier Epoch~5 F200W mosaic from the NEXUS Quick Data Release \citep{Zhuang2026NEXUSQDR}, drizzled to $0.03''$\,pixel$^{-1}$. The region we use spans about $12.6'\times11.8'$ ($\mathrm{RA}=268.21^\circ$--$268.71^\circ$, $\mathrm{Dec}=65.10^\circ$--$65.30^\circ$), of which $\sim\!70\,\mathrm{arcmin}^2$ has F200W coverage. 
At $\sim\!0.07''$, the F200W PSF is about 2.5 times narrower than the $I_{\rm E}$ mosaic PSF ($\sim\!0\farcs17$) and spans only $\sim\!1.4$ pixels of our $0.05''$ reconstructions. 
Because NEXUS and Q1 observe the same galaxies, F200W lets us check whether structures recovered from real Euclid data appear in a sharper image of the
same objects. 
It does not, however, provide photometric truth: at $\sim\!2\,\mu$m the
light traces different stellar populations and suffers less dust attenuation than in $I_{\rm E}$, so the comparisons are mostly qualitative.

\section{Methods}
\subsection{Training data}
\subsubsection{Measuring PSF}
\subsubsection{Measuring PSF}
To simulate Euclid observations of a synthetic sky we first need the Euclid point-spread function (PSF).
The PSF characterizes how Euclid blurs point sources, influenced by diffraction (including the spikes of the VIS secondary mirror spider mount), optical aberrations, detector effects and pointing instability.
Rather than adopting an instrument model, we measure the effective PSF (ePSF; \citealt{AndersonKing2000}) directly from bright point sources in the Q1 MER mosaics with \texttt{photutils}.
Importantly, the ePSF also captures dithering and resampling, which is significant for NISP's $0.3''$ native pixels resampled to $0.1''$ pixels by the MER pipeline.
We select $\sim$43{,}000 unmasked point-like sources with $16 < \mathrm{VIS} < 19$ from the MER catalogue, extract $51''$ cutouts in all four bands, and reject sources with saturated cores, keeping 17{,}905 stars usable in every band.
To capture spatial variation we cluster them by sky position into 44 regions of $\sim$400 stars each and build one ePSF per region and band, oversampled by a factor of two ($0.05''$ pixels).
Finally, we multiply each ePSF stamp by a radial cosine window that brings it smoothly to zero towards the stamp edge, so the kernel has no discontinuity at its boundary.
The resulting median FWHMs are $0.15''$ (VIS), $0.36''$ ($Y_E$), $0.38''$ ($J_E$) and $0.39''$ ($H_E$).

Because the mosaics are north-aligned while the PSF is fixed to the instrument, the PSF orientation on the sky depends on the telescope roll angle.
To account for that, we augment the ePSFs with random rotations.
We also apply a small random elastic warp to each kernel, following \citet{Connor2022}, which perturbs the core width (by $\approx 10\%$) and the faint structure while conserving flux.
This keeps the network from overfitting to our finite set of measured ePSFs, which carry extraction errors and do not resolve all of the PSF variation across the field.


\begin{figure*}
    \centering
    \includegraphics[width=1\linewidth]{images/psfs_idx0_VIS-Y_E_VIS_figure.png}
    \caption{measuring the psfs}
    \label{fig:psfs}
\end{figure*}
        
\subsubsection{Noise \& Artifacts}
After convolving the synthetic sky with the PSF, we downsample it by a factor of two to the $0.1''$ MER pixel scale.
Next, we simulate the noise.
Each scene draws its noise level from the Q1 MER RMS maps, measured at 294 random extragalactic positions, so the training set spans the real range of survey depths (median RMS of 29.8, 9.4, 9.9 and 9.7\,e$^-$ per pixel in VIS, $Y_E$, $J_E$ and $H_E$).
On top of the background we add Poisson noise from the sources themselves.
Because the MER mosaics are built by resampling several dithered exposures, their noise is correlated between neighbouring pixels, most strongly in NISP, where $0.3''$ native pixels are resampled to $0.1''$.
We reproduce this by generating white noise for four exposures on the native detector grid, shifting each by a random sub-pixel dither, resampling it to the mosaic grid and averaging.
We also vary the depth slightly across each field and occasionally add an edge where the depth steps between two pointings, as seen in the real mosaics.

Finally, we add the artifacts that survive MER processing.
Most cosmic rays and hot pixels are rejected when the exposures are stacked, so we inject only the small fraction that remains, together with dead pixels and faint streaks that commonly confuse the model.
We also add straight diffraction spikes across the fields, emulating bright stars far away from the field. 
Saturation is modelled as in MER, which masks saturated cores by zeroing pixels with a probability that rises with how far they exceed the full-well level.
\begin{figure*}
    \centering
    \includegraphics[width=1\linewidth]{PSF_Noise/noise_artifacts.png}
    \caption{Noise and Artifacts}
    \label{fig:noiaseandartifacts}
\end{figure*}
\subsubsection{Distributions \& Transformations}
We first set the number, brightness and size of the galaxies in each training scene. 
These distributions must match Euclid's, because any difference would be learned by the network and reflected in the reconstructions of real data.

We start my modelling brightness counts.
The measure of brightness that MER catalogue provides is VIS flux in 2FWHM aperture, where FWHM is taken of an external ground band, which is about $1.3''$, or $13$ pixels in Euclid VIS resolution.
We find that galaxy counts grow smoothly toward faint magnitudes with almost-linear density/magnitude law, until Q1 data becomes incomplete near VIS $\approx 25$. 
We fit this growth with a segmented-linear relation and augment with flat extrapolation up to VIS $=29$.

Each galaxy also needs a size, which we draw conditional on its brightness.
For a sample of $\approx7.5$ million Q1 galaxies whose Sérsic fits passed MER's quality checks, we take the circularized half-light radius $R_e=R_{e,\mathrm{major}}\sqrt{q}$ as the size, where $R_{e,\mathrm{major}}$ is the half-light radius along the major axis and $q$ the axis ratio
We model the distribution of $R_e$ at a fixed magnitude as log-normal, with a median that decreases as a power law toward fainter magnitudes, and width that is the same at all magnitudes.
We fit these three parameters to the found and binned Q1 radii between VIS magnitudes $17$ and $25.5$, capping the weight of each magnitude bin so that the numerous faint galaxies do not outweight the rare bright one.

Finally, each galaxy needs colors, which set its brightness in the NISP bands relative to VIS, and a star-formation rate (SFR), which we will use to select the TNG galaxy with similar morphology.
At fixed brightness and size these quantities have a complex correlated distribution, so we use a non-parametric adaptive model: a random forest.
It is trained on approximately $8.4\times 10^4$ Q1 galaxies and divides the magnitude--radius space into regions of similar colors and SFR.
We use extreme deconvolution method \citep{Bovy2011} to remove the measurement noise from each region.
\begin{figure*}
    \centering
    \includegraphics[width=1\linewidth]{images/euclidpolish_galaxy_distributions_2x2.png}
    \caption{Galaxy population}
    \label{fig:galaxieees}
\end{figure*}

Stars constitute the second population in our scenes. 
Resolving them is not among our goals, and they rather act as artifacts, bringing the PSF, spikes and saturated cores into the scenes. 
We collect the statistics of density of star-like objects, as reported by PHZ catalogue, binned by AB magnitude in VIS band (Figure~\ref{fig:star_population}a).
This results in an almost-linear law in log density, that falls for objects with magnitude lower than $18$ and higher than $21$.
We fit a straight line to the linear part, extrapolating a simple magnitude-density dependence.
As a check, we compare it to Gaia DR3 (even though it's $G$ band differs from Euclid's VIS band): its counts agree with the line to $\sim\!10\%$ for $16 < G < 21$, and drop below it only for the few stars brighter than $G\approx 14$.

Each synthetic star also needs fluxes in the three NISP bands, which we set through its colors VIS$-Y_{\rm E}$, $Y_{\rm E} - J_{\rm E}$ and $J_{\rm E} - H_{\rm E}$.
We use an approach similar, to how we modeled galactic distributions. 
We therefore resample real stars, as we do for galaxies, but with brightness as the only condition: a synthetic star takes the colors of a Q1 star from the same $0.5$-mag bin, picked in proportion to the spectroscopic (PHZ) star probability.
Within each bin, extreme deconvolution \citep{Bovy2011} of the three colors jointly removes the measurement noise, so the copies carry the intrinsic variance rather than the catalog errors.
This follows the reddenning trend of stars towards faint magnitudes and, once the Q1 errors are added back, reproduces the observed color widths to within $7\%$ (Figure~\ref{fig:star_population}b--d).
Positions of stars are random and are drawn on the noise-stage of data (so each galactic field can get different stars after adding noise with different seeds).
\begin{figure*}
    \centering
    \includegraphics[width=\textwidth]{images/star_population.pdf}
    \caption{Star population of the training scenes}
    \label{fig:star_population}
\end{figure*}
\subsubsection{Resulting pairs}
\begin{figure*}
    \centering
    \includegraphics[width=1\linewidth]{images/fig04_synthetic_lr_sr_hr_fields.png}
    \caption{Synthetic triplets}
    \label{fig:synthetictriples}
\end{figure*}
\subsection{Model}
Our model is a residual convolution network based on the WDSR \citep{Yu2018} and extending POLISH \citep{Connor2022}.
We keep it small, since a survey-wide reconstruction needs a model that is cheap to run.
It takes the four MER bands on the $0.1''$ grid and returns all four on a $0.05''$ grid (Figure ~\ref{fig:architecture}).
Two paths run side by side: \begin{itemize}
    \item The main path mixes all bands from its first layer, passes them through $16$ residual blocks of $32$ feature maps.
    \item The second path is a light upsamples for each band that sees only that band, and the main path adds its correction to it. 
\end{itemize}
With about $3\times 10^5$ parameters and no fully connected layers, the model runs within minutes on large mosaics.

The light in our scenes covers a range of $6$ orders of magnitude.
POLISH++ compresses such a range with a single $\mathrm{asinh}$ \citep{Wu2026}, but one choice is left: what should be the units in which we pass the images through $\mathrm{asinh}$? 
Since this operator compresses logarithmically all the values $\gg 1$ and leaves approximately linear values $\approx 0$, it will act differently depending on how we scale the images first.
Difference of our approach from POLISH++ is in choice of the ``knees``: the normalization constants $$\mathrm{asinh}(\text{pixel} / \text{knee}).$$
Instead of choosing one, our loss and input use eight: knees of $10^{-3}$, $10^{-2}$, \ldots, $10^{4}$ electrons, covering the full range of usual input data.
This replaces the single stretch of POLISH++, which forces a single choice for a dynamic range that spans more than seven decades in our images. 
Since the data inside of the model is transformed with $\mathrm{asinh}$ we should also add a $\text{knee}\cdot\mathrm{sinh}$ operation before the output. 
The knee for this operation is trained as one of the parameters of the model, optimized for the best output.

Training minimizes a sum of two terms.
The first compares reconstruction and target after stretching

\task{Architecture figure: work more on the labels so people can read it immediately, with no thinking.}

% Architecture figure: the whole figure* block (tikzpicture + caption).
% Needs architecture_preamble.tex (colours, styles, \mdfmap/\mdslabs/\mdimstack,
% the \mdfun pill and the \mdg* curve glyphs) and images/thumbs/{lr,sr}_rgb.png.
% Conventions: cuboids are feature maps (gray = 0.1'' grid, green = 0.05'' grid,
% number above = channels); arrows carry the linear steps; the purple pills
% are the only pixel-wise non-linear functions.
% No blank lines inside the tikzpicture: with the linenumbers class option each
% one would end the build with "Improper \prevdepth".
\begin{figure*}
\centering
\begin{tikzpicture}
% ===================================================================== (a)
\def\hL{1.45}   % 0.1'' grid
\def\hH{2.35}   % 0.05'' grid
\def\ys{3.35}   % per-band row
\def\lrs{1.25}  % input image side; the output image is exactly twice this
\def\yi{-2.60}  % centre of the input image
\def\yo{-3.30}  % centre of the output image
\def\xs{15.95}  % x of the sum and of the output column
\node[mdpanel] at (0,6.40) {(a)};
% ---- legend
\begin{scope}[shift={(0.90,6.20)}]
  \mdfmap{lg1}{0}{0}{0.22}{0.42}{mdLR}
  \node[mdleg,anchor=west] at (0.45,0) {feature maps, $0.1''$ grid};
  \mdfmap{lg2}{4.10}{0}{0.22}{0.42}{mdHR}
  \node[mdleg,anchor=west] at (4.55,0) {feature maps, $0.05''$ grid};
  \node[mdch,anchor=west] at (8.45,0) {32};
  \node[mdleg,anchor=west] at (8.80,0) {number of channels};
  \node[mdfun,font=\fontsize{8}{9.5}\selectfont,inner sep=2pt] at (12.20,0) {\mdgrelu};
  \node[mdleg,anchor=west] at (12.65,0) {pixel-wise non-linear function};
\end{scope}
% ---- input, stretch, split
\mdimstack{in}{0}{\yi}{\lrs}{lr_rgb}
\node[mdop,anchor=north west,align=left] at ($(in-s)+(-0.5*\lrs,-0.06)$)
  {Euclid image $x$\\[-1pt] 4 bands, $0.1''$ pixels};
\coordinate (riseI) at ($(in-s)+(0.105,\lrs+0.21)$);
\node[mdfun] (asinh) at ($(riseI)+(0,1.30)$) {\mdgasinh\ \mdtwo{$\mathrm{arcsinh}(x/k_j)$}{8 knees $k_j$}};
\draw[mdflow] (riseI) -- (asinh.south);
\coordinate (split) at (asinh.north |- 0,0);
\fill[mdLine] (split) circle (1.6pt);
\draw[mdline] (asinh.north) -- (split);
% ---- shared trunk
\mdfmap{S}{2.40}{0}{0.72}{\hL}{mdLR}
\mdslabs{2.40}{0}{0.72}{\hL}{8}{mdLR}
\draw[mdflow] (split) -- (S-w);
\node[mdch,above=1pt of S-n] {$4\times8 = 32$};
\mdfmap{F}{5.60}{0}{0.72}{\hL}{mdLR}
\draw[mdflow] (S-e) -- (F-w) node[mdop,midway,above=1pt] {$3\!\times\!3$\\ convolution};
\node[mdch,above=1pt of F-n] {32};
% 16 residual blocks: a stack of three boxes, the front one labelled
\foreach \k in {2,1}{%
  \draw[rounded corners=2pt,line width=0.6pt,draw=mdRes!75!black,fill=mdRes!22]
    ($(7.27,-0.5*\hL)+(0.11*\k,0.11*\k)$) rectangle ++(1.90,\hL);}
\node[mdbox,draw=mdRes!75!black,fill=mdRes!38,minimum width=1.90cm,minimum height=\hL cm]
  (res) at (8.22,0) {16\\ residual\\ blocks};
\draw[mdflow] (F-e) -- (res.west);
\mdfmap{T}{11.32}{0}{0.59}{\hL}{mdLR}
\draw[mdflow] (res.east) -- (T-w) node[mdop,midway,above=1pt] {$3\!\times\!3$\\ convolution};
\node[mdch,above=1pt of T-n] {16};
\mdfmap{M}{14.42}{0}{0.33}{\hH}{mdHR}
\draw[mdflow] (T-e) -- (M-w) node[mdop,midway,above=1pt] {upsample $2\times$\\ (pixel shuf\/f\/le)};
\node[mdch,above=1pt of M-n] {4};
\node[mdnote,anchor=north west] at (2.40,-0.98) {shared trunk: all bands mixed};
% ---- sum, learned knee, output (exactly twice the input image)
\node[mdsum] (sum) at (\xs,0) {$+$};
\draw[mdflow] (M-e) -- (sum);
\node[mdfun] (knee) at (\xs,-1.12) {\mdgsinh\ \mdtwo{$k_b\sinh y_b$}{learned knee $k_b$}};
\draw[mdflow] (sum) -- (knee.north) node[mdop,midway,right=1pt] {$y_b$};
\mdimstack{out}{\xs-\lrs-0.105}{\yo}{2*\lrs}{sr_rgb}
\draw[mdflow] (knee.south) -- ($(out-s)+(0.105,2*\lrs+0.21)$);
\node[mdop,anchor=north,align=center] at ($(out-s)+(0,-0.06)$)
  {reconstruction\\[-1pt] 4 bands\\[-1pt] $0.05''$ pixels};
% ---- per-band path (one copy per band, drawn as a stack of four sheets)
\foreach \k in {3,2,1}{%
  \draw[rounded corners=3pt,line width=0.5pt,draw=black!30,fill=white]
    ($(4.55,1.95)+(0.13*\k,0.13*\k)$) rectangle ++(7.15,3.45);}
\draw[rounded corners=3pt,line width=0.6pt,draw=black!55,dashed,fill=white]
  (4.55,1.95) rectangle ++(7.15,3.45);
\node[mdnote,anchor=south west] at (4.68,2.03) {each band $b$ separately (all 4 in parallel)};
\mdfmap{K}{5.00}{\ys}{0.46}{\hL}{mdLR}
\mdslabs{5.00}{\ys}{0.46}{\hL}{8}{mdLR}
\mdfmap{KC}{7.75}{\ys}{0.33}{\hL}{mdLR}
\mdfmap{KS}{10.65}{\ys}{0.07}{\hH}{mdHR}
\mdfmap{KA}{14.42}{\ys}{0.33}{\hH}{mdHR}
\draw[mdflow,rounded corners=5pt] (split) |- (K-w)
  node[mdop,pos=0.75,above=1pt] {band $b$ only};
\node[mdch,above=1pt of K-n] {8};
\draw[mdflow] (K-e) -- (KC-w) node[mdop,midway,above=1pt] {$5\!\times\!5$\\ convolution};
\node[mdch,above=1pt of KC-n] {4};
\draw[mdflow] (KC-e) -- (KS-w) node[mdop,midway,above=1pt] {upsample $2\times$\\ (pixel shuf\/f\/le)};
\node[mdch,above=1pt of KS-n] {1};
\draw[mdflow] (KS-e) -- (KA-w) node[mdop,midway,above=1pt,xshift=4pt] {stack the\\ 4 bands};
\node[mdch,above=1pt of KA-n] {4};
\draw[mdflow,rounded corners=5pt] (KA-e) -| (sum.north);
% ===================================================================== (b)
% same notation as (a); 192 -> 1.06, 32 -> .72, 25 -> .67
\def\yb{-6.05}
\def\hb{1.30}
\def\bT{-4.35}  % top of the (b) frame
\def\bB{-7.45}  % bottom of the (b) frame
\node[mdpanel] at (0,\bT+0.02) {(b)};
% zoom wedge from the residual blocks in (a) to the frame of (b)
\begin{scope}[on background layer]
  \fill[mdRes!12] (res.south west) -- (res.south east) -- (11.6,\bT) -- (4.2,\bT) -- cycle;
\end{scope}
\draw[mdRes!70!black,line width=0.5pt] (res.south west) -- (4.2,\bT) (res.south east) -- (11.6,\bT);
\draw[rounded corners=4pt,draw=mdRes!85!black,line width=0.9pt,fill=white] (0.62,\bB) rectangle (14.50,\bT);
\node[font=\fontsize{9}{10.5}\selectfont\bfseries,text=mdRes!65!black,anchor=north west,inner sep=0pt]
  at (0.80,\bT-0.14) {one residual block \mdseries(16 in a row in panel a)};
\mdfmap{bx}{1.00}{\yb}{0.72}{\hb}{mdLR}
\mdfmap{b1}{4.05}{\yb}{1.06}{\hb}{mdLR}
\node[mdfun] (relu) at (6.55,\yb) {\mdgrelu\ ReLU};
\mdfmap{b2}{9.15}{\yb}{0.67}{\hb}{mdLR}
\mdfmap{b3}{12.10}{\yb}{0.72}{\hb}{mdLR}
\node[mdsum] (bsum) at (13.78,\yb) {$+$};
\draw[mdflow] (bx-e) -- (b1-w) node[mdop,midway,above=1pt] {$1\!\times\!1$\\ convolution};
\draw[mdflow] (b1-e) -- (relu.west);
\draw[mdflow] (relu.east) -- (b2-w) node[mdop,midway,above=1pt] {$1\!\times\!1$\\ convolution};
\draw[mdflow] (b2-e) -- (b3-w) node[mdop,midway,above=1pt] {$3\!\times\!3$\\ convolution};
\draw[mdflow] (b3-e) -- (bsum);
\draw[mdflow] (bsum) -- (14.50,\yb);
\draw[mdflow,rounded corners=5pt] ($(bx-s)+(0.06,0)$) -- ++(0,-0.48) -| (bsum.south)
  node[mdop,pos=0.25,above=1pt] {block input, added back unchanged};
\node[mdch,above=1pt of bx-n] {32};
\node[mdch,above=1pt of b1-n] {192};
\node[mdch,above=1pt of b2-n] {25};
\node[mdch,above=1pt of b3-n] {32};
\end{tikzpicture}
\caption{The reconstruction network.
(a) The four-band Euclid image $x$ (electrons, $0.1''$ pixels) is stretched with $\mathrm{arcsinh}(x/k_j)$ at eight knees $k_j = 10^{-3}, 10^{-2}, \ldots, 10^{4}$ electrons, giving $4\times8=32$ channels.
Two paths start from this stack.
The shared trunk (bottom row) mixes all bands: a $3\times3$ convolution, $16$ residual blocks, a second $3\times3$ convolution and a pixel shuffle, which rearranges each group of four channels into a $2\times2$ patch of pixels on the $0.05''$ grid.
The per-band path (top row) upsamples each band from its own eight channels with a $5\times5$ convolution and a pixel shuffle, so it never mixes bands.
The two are added, giving the stretched image $y_b$ for each band $b$, and $k_b \sinh y_b$ with a learned knee $k_b$ turns it back into electrons.
Gray cuboids are feature maps on the $0.1''$ grid and green ones on the $0.05''$ grid; numbers above them are channel counts.
Purple boxes are the only pixel-wise non-linear functions; every other step is linear.
(b) One of the $16$ residual blocks: a $1\times1$ convolution widens $32$ channels to $192$, a ReLU, a $1\times1$ convolution narrows them to $25$, and a $3\times3$ convolution returns $32$ channels, which are added to the block input.
The thumbnails show a Q1 lens candidate and its reconstruction (R\,=\,$H_{\rm E}$, G\,=\,mean of $Y_{\rm E}$ and $J_{\rm E}$, B\,=\,VIS).}
\label{fig:architecture}
\end{figure*}


\section{Results}
\subsection{Metrics}
A table of models + combiner with their metrics.
Plots (cross-correlation, error-std).
\begin{figure*}
    \centering
    \includegraphics[width=1\linewidth]{images/fig09_ensemble_diagnostics.png}
    \caption{METRICS}
    \label{fig:metrics}
\end{figure*}
\subsection{Images}
Synthetic images +  real images + lenses.
\begin{figure*}
    \centering
    \includegraphics[width=1\linewidth]{images/fig07_nexus_closeup_comparisons.png}
    \caption{REAL imagEs}
    \label{fig:realImAges}
\end{figure*}
\begin{figure*}
    \centering
    \includegraphics[width=1\linewidth]{images/fig04_synthetic_lr_sr_hr_fields.png}
    \caption{Triplets AGAIN}
    \label{fig:3plets}
\end{figure*}

\section{Discussion}
\begin{itemize}
    \item Compare to POLISH (ensemble! TNGs!)
    \item Compare to other SR papers
    \item Discussion of different architectures
    \item Future directions, including adaptability to other telescopes (Roman)
\end{itemize}
\section{Conclusion}
Summary

\bibliography{references}{}
\bibliographystyle{aasjournalv7}


\end{document}

% End of file `sample7.tex'.
